\documentclass{beamer}
\usepackage[utf8]{inputenc}
\usetheme{Berlin}
\usecolortheme{beaver}
\setbeamertemplate{page number in head/foot}[totalframenumber]
\setbeamercovered{invisible}
\setbeamercovered{again covered={\opaqueness<1->{30}}}
\usepackage{multicol}
\usepackage{ dsfont }
\usepackage{listings}
\usepackage{fancyvrb}
\usepackage{verbatimbox}
\usepackage{tikz}
\usepackage{array}
\usepackage{amsmath}
\usepackage{geometry}
\usepackage{longtable}
\usepackage{amsthm}
\usepackage{amssymb}
\usepackage{appendix}
\usepackage{dsfont}
\usepackage{tikz}
\usepackage{tikz-qtree}
\usepackage{mathrsfs}
\usepackage{latexsym}
\usepackage{caption}
\usepackage{subcaption}
\usepackage{multicol}
\usepackage{xcolor}
\newcommand\Fontvi{\fontsize{9}{7.2}\selectfont}
\usepackage{cancel}


\usepackage{amsmath, amsfonts, amscd, amssymb, amsthm, graphicx, stmaryrd}
\usepackage[all,2cell,ps]{xy}




\title[Understanding Musical Chords Through Geometric Transformations] %optional
{Understanding Musical Chords Through Geometric Transformations}


\author % (optional, for multiple authors)
{Steven Kesser \& Elizabeth McArthur}

\institute[] % (optional)
{
  University of Georgia \\ Advised By: \\Cameron Thomas 
}

\date[April 28, 2022] % (optional)
{Directed Reading Program\\ November 22, 2024}

\logo

\begin{document}

\frame{\titlepage}

\begin{frame}
\frametitle{Table of Contents}
\tableofcontents
\end{frame}

\section{Motivation}

\begin{frame}{Chord Progressions}
\begin{block}{Definition}
     A \textit{chord progression} is the movement from one chord to the next.
\end{block}

\begin{itemize}
    \item Most musical styles are largely determined by the chord progressions used in the music.
    \item The most common and strongest chord progressions are circle progressions.
\end{itemize}

\begin{Example}
    The following is a ii-V-I circle progression in the key of C:
    \begin{center}
        Cmaj - Dmin - Gmaj - Cmaj
    \end{center}
\end{Example}
\end{frame}

\begin{frame}{Motivating Question}

\begin{alertblock}{Question}
     Let $\mathcal{M}$ the space of all triads. Can we classify chord progressions by classifying paths in $\mathcal{M}$? More specifically, Can we classify circle progressions by understanding maps $S^1 \longrightarrow \mathcal{M}$ (aka fundamental group $\pi(\mathcal{M})$)?
\end{alertblock}
\end{frame}

\section{Music Theory}
\begin{frame}{Pitch}
    \begin{block}{Definition}
     \textit{Pitch} is the highness or lowness of a sound. Variations in frequency are what we hear as
     variations in pitch. 
    \end{block}
 \begin{Definition} 
A pitch class contains all the same pitches regardless of octave.
 \end{Definition}

\end{frame}

\begin{frame}{Scales}
    \begin{block}{Definition}
    Scales are an collection of pitches in ascending and descending order that differ by half(H) and whole(W) steps.
    \end{block}
    
    \begin{Example}
    A major scale is constructed by a series below:
    \begin{center}
    W W H W W W H 
    \end{center}
    This can be seen from the construction of the C major scale from a chromatic scale, a scale with all 12 unique pitches.
    \begin{center} \textcolor{blue}{\textbf{C}} C$\sharp$ \textcolor{blue}{\textbf{D}} D$\sharp$ \textcolor{blue}{\textbf{E}} \textcolor{blue}{\textbf{F}} F$\sharp$ \textcolor{blue}{\textbf{G}} G$\sharp$ \textcolor{blue}{\textbf{A}} A$\sharp$ B \textcolor{blue}{\textbf{C}}
\end{center}
    \end{Example}
\end{frame}
\begin{frame}{Scales}
\begin{Example}
    A minor scale is constructed by a series of
    \begin{center}
    W H W W H W W
    \end{center}
    This can be seen from the construction of the C minor scale from a C minor chromatic scale, a scale with all 12 unique pitches.
    \begin{center} \textcolor{blue}{\textbf{C}} D$\flat$
    \textcolor{blue}{\textbf{D}} \textcolor{blue}{\textbf{E$\flat$}} E \textcolor{blue}{\textbf{F}} G$\flat$ \textcolor{blue}{\textbf{G}} \textcolor{blue}{\textbf{A$\flat$}} A \textcolor{blue}{\textbf{B$\flat$}} B \textcolor{blue}{\textbf{C}}
\end{center}
    \end{Example}
\end{frame}

\begin{frame}{Harmony}

\begin{Definition}
   Harmony is the musical result of tones sounding together, typically expressed by chords.
\end{Definition}
\begin{Definition}
A chord is a harmonic unit with at least three different tones sounding simultaneously. 

\end{Definition}
\end{frame}



\begin{frame}{Triads}
    \begin{block}{Definition}
    A triad is a three-toned chord.\\
    \begin{itemize}
        \item Triads typically come in four qualities: major, minor, diminished and augmented.
    \end{itemize}
    \end{block}
    
    \begin{center}
    \includegraphics[width=75mm]{C-Major-Piano-Chord-Illustration-showing-keyboard-and-treble-staff-notes-C-E-G}
    \end{center}
   
\end{frame}

\begin{frame}{Major Triad Example}
 \begin{Example}
    A major triad consists of a the following intervals:
    \begin{itemize}
        \item Root - The first note in the scale
        \item Major Third - fourth half steps from root
        \item Perfect Fifth - seven half steps from root
    \end{itemize}

    Given the C major scale, we can construct a C major triad(in red):
    \begin{center} \textcolor{red}{\textbf{C}} C$\sharp$ \textcolor{blue}{\textbf{D}} D$\sharp$ \textcolor{red}{\textbf{E}} \textcolor{blue}{\textbf{F}} F$\sharp$ \textcolor{red}{\textbf{G}} G$\sharp$ \textcolor{blue}{\textbf{A}} A$\sharp$ B \textcolor{blue}{\textbf{C}}
    \end{center}
\end{Example}
\end{frame}

\begin{frame}{Minor Triad Example}
 \begin{Example}
    A minor triad consists of a the following intervals:
    \begin{itemize}
        \item Root - The first note in the scale
        \item Minor Third - three half steps from root
        \item Perfect Fifth - seven half steps from root
    \end{itemize}

    Given the C major scale, we can construct a C major triad(in red):
    \begin{center}
    \textcolor{red}{\textbf{C}} D$\flat$
    \textcolor{blue}{\textbf{D}} \textcolor{red}{\textbf{E$\flat$}} E \textcolor{blue}{\textbf{F}} G$\flat$ \textcolor{red}{\textbf{G}} \textcolor{blue}{\textbf{A$\flat$}} A \textcolor{blue}{\textbf{B$\flat$}} B \textcolor{blue}{\textbf{C}}
    \end{center}
\end{Example}
\end{frame}

\section{Preserving Chord Quality}

\begin{frame}{Triads as Triangles}
    We can encoded the 12 unique pitches from the chromatic scale on the unit circle depicted below: 
    \begin{center}
    \includegraphics[width=42mm]{Screenshot 2024-11-06 213028}
    \end{center}
    Remark: We represent each pitch with the point $(\cos(\frac{i\pi}{6}),\sin(\frac{i\pi}{6}))$
    \end{frame}

\begin{frame}{Intervals}
\begin{itemize}
    \item We can also construct tridas based on semitonal intervals. Start with a root and choose two semitonal intervals. The choice of semitonal intervals gives the chord its quality.
\end{itemize}
\begin{Example}
\begin{multicols}{2}
\begin{itemize}
    \item minor third  (m3) = three semitones above root (-3)
    \item major third  (M3) = four semitones above root (-4)
    \item perfect fourth  (P4) = five semitones above root (-5)
    \item tritone (Dim4) = six semitones above root (-6)
    \item perfect fifth (P5) = seven semitones above root (-7)
    \item minor sixth (m6) = eight semitones above root (-8)
\end{itemize}
\end{multicols}
\end{Example}


\end{frame}

\begin{frame}{Triads as Triangles}
    We can construct a 2x3 matrix that is able to encode all major triads.\\
    \vspace{3mm}
    Let k be the root, and $I_1$ and $I_2$ be the two semitonal intervals in a given triad. Then the follow matrix encodes the given triad.
    \[
    k_{\text{Q}} = 
    \begin{pmatrix}
      \cos\left(\frac{k\pi}{6}\right) & \cos\left(\frac{(k+I_1)\pi}{6}\right) & \cos\left(\frac{(k+I_1)\pi}{6}\right) \\
      \sin\left(\frac{k\pi}{6}\right) & \sin\left(\frac{(k+I_2)\pi}{6}\right) & \sin\left(\frac{(k+I_2)\pi}{6}\right)
    \end{pmatrix}
    \]
\end{frame}

\begin{frame}{Example of Triads as Triangles}
    Recall, a major triad consists of a the following intervals:
    \begin{itemize}
        \item Root - The first note in the scale
        \item Major Third - fourth half steps from root
        \item Perfect Fifth - seven half steps from root
    \end{itemize}

    Thus the follow matrix encodes a major chord.

    \[
    k_{\text{maj}} = 
    \begin{pmatrix}
      \cos\left(\frac{k\pi}{6}\right) & \cos\left(\frac{(k-4)\pi}{6}\right) & \cos\left(\frac{(k-7)\pi}{6}\right) \\
      \sin\left(\frac{k\pi}{6}\right) & \sin\left(\frac{(k-4)\pi}{6}\right) & \sin\left(\frac{(k-7)\pi}{6}\right)
    \end{pmatrix}
    \]

   Remark: k corresponds to the root of the triad
\end{frame}
    
\begin{frame}{Example of Triads as Triangles}
    The matrix corresponding to a C major triad is
     \[
    3_{\text{maj}} = 
    \begin{pmatrix}
      \cos\left(\frac{3\pi}{6}\right) & \cos\left(\frac{-\pi}{6}\right) & \cos\left(\frac{(-4)\pi}{6}\right) \\
      \sin\left(\frac{3\pi}{6}\right) & \sin\left(\frac{-\pi}{6}\right) & \sin\left(\frac{(-4)\pi}{6}\right)
    \end{pmatrix}
    \]
    \begin{center}
        \includegraphics[width=40mm]{Screenshot 2024-11-07 095805.png}
    \end{center}
    
\end{frame}



\begin{frame}
{Rotation Matrix}
\begin{Definition}
A rotation matrix is an orthogonal matrix that rotates vectors in Euclidean space.
\begin{center}
$R_{\theta} = \begin{pmatrix}
  \cos{\theta} & -\sin{\theta}\\ 
  \sin{\theta} & \cos{\theta}
\end{pmatrix}$
\end{center}
\end{Definition}
\end{frame}


\begin{frame}{Rotation Preserves Chord Quality}
We act on our triangle with a rotation matrix where $\theta = \frac{i\pi}{6}$ by $R_{\theta}*k_{Q} = R_{\theta}k_{Q}$

\begin{alertblock}{Theorem}
    Let $R_{\theta}$ be the rotation matrix by $\theta$  where $\theta = \frac{i\pi}{6}$ for $0 \leq i \leq 11$. Also, let $k_{Q}$ be the matrix for the k triad for a given quality. Then, 
    \begin{center}
        $R_{\theta} \cdot k_{Q} = (k+i)_{Q}$ 
    \end{center}
\end{alertblock}
\end{frame}

\begin{frame}{Sketch of Proof}
When compute the product $R_{\theta} \cdot k_{Q}$, our resultant matrix is as follows:
{\small \[\begin{pmatrix}
  \cos\left(\frac{(i+k)\pi}{6}\right) & \cos\left(\frac{(i+k+I_1)\pi}{6}\right) & \cos\left(\frac{(i+k+I_2)\pi}{6}\right) \\
  \sin\left(\frac{(i+k)\pi}{6}\right) & \sin\left(\frac{(i+k+I_1)\pi}{6}\right) & \sin\left(\frac{(i+k+I_2)\pi}{6}\right)
\end{pmatrix}
\]}

Which we can see takes the form of our already defined major matrix. Thus rotation preserves quality of chord.
{\small  \[
\begin{pmatrix}
  \cos\left(\frac{k\pi}{6}\right) & \cos\left(\frac{(k+I_1)\pi}{6}\right) & \cos\left(\frac{(k+I_2)\pi}{6}\right) \\
  \sin\left(\frac{k\pi}{6}\right) & \sin\left(\frac{(k+I_1)\pi}{6}\right) & \sin\left(\frac{(k+I_2)\pi}{6}\right)
\end{pmatrix}
\]}

Remark: This calculation works out for any chord quality.

\end{frame}

\begin{frame}{Visualization of Rotations}
    \includegraphics[width=44mm]{Screenshot 2024-11-07 095805.png}
    \hspace{5mm} 
    \raisebox{20mm}{$\rightarrow$} 
    \hspace{5mm} 
    \includegraphics[width=40mm]{Screenshot 2024-11-07 095724.png}\\
     \vspace{4mm}
    Remark: This corresponds with a rotation of $5\pi/6$.
\end{frame}

\begin{frame}{Visualization of Rotations}
    \includegraphics[width=44mm]{Screenshot 2024-11-07 095805.png}
    \hspace{5mm} 
    \raisebox{20mm}{$\cancel{\rightarrow}$} 
    \hspace{5mm} 
    \includegraphics[width=40mm]{Screenshot 2024-11-07 095856.png}\\
    \vspace{4mm}
    Remark: There is no orthogonal transformation that can change the quality of chord. Thus we need to find a map that alters either the second or thrid column of our chordal matrices.
    
\end{frame}

\begin{frame}{More Results}
Augmented chords form equilateral triangles. Thus, their quality is invariant under the three axes of reflectional symmetry.

     \begin{center}
        \includegraphics[width=31mm]{Screenshot 2024-11-19 183927.png}
        \includegraphics[width=31mm]{Screenshot 2024-11-19 191137.png}
    \end{center}
    Diminished chords form isosceles triangles. Thus, their quality is invariant under one axis of reflection.
\end{frame}

\begin{frame}{Project Prospects}

\begin{itemize}
    \item Currently, we are exploring the various transformations on the triad and how they effect the quality of the chord.
    \item In later research we hope to be able to detect chord progressions in songs, in hopes of finding patterns in relation to genre, emotion, and popularity.
\end{itemize}
  
\end{frame}




\section{Conclusion}

\begin{frame}{References}
\begin{thebibliography}{99}

Benward, Bruce, and Marilyn Saker. Music in Theory and Practice. 8th ed., vol. 1, McGraw-Hill Education, 2020.


\end{thebibliography}
\end{frame}


\begin{frame}{Thank You!}

\centering
\Huge Questions?
    
\end{frame}
\end{document}

\begin{block}{Rotating $v_1$}
\begin{center}
$R_{\frac{i\pi}{6}}v_{1} = \begin{pmatrix}
  \cos{\frac{i\pi}{6}} & -\sin{\frac{i\pi}{6}}\\ 
  \sin{\frac{i\pi}{6}} & \cos{\frac{i\pi}{6}}
\end{pmatrix} \begin{pmatrix}
  \cos({\frac{k\pi}{6}}) \\ 
  \sin({\frac{k\pi}{6}})
\end{pmatrix} = \begin{pmatrix}
  \cos({\frac{i\pi}{6}})\cos({\frac{k\pi}{6}}) - \sin({\frac{i\pi}{6}})\sin({\frac{k\pi}{6}})\\ 
  \sin({\frac{i\pi}{6}})\sin({\frac{k\pi}{6}}) + \cos({\frac{i\pi}{6}})\cos({\frac{k\pi}{6}})
\end{pmatrix} = \begin{pmatrix}
  \cos({\frac{(i+k)\pi}{6}}) \\ 
  \sin({\frac{(i+k)\pi}{6}})
\end{pmatrix}$
\end{center}
\end{block}

\begin{frame}
\begin{block}{Rotating $v_2$}
\begin{center}
$R_{\frac{i\pi}{6}}v_{2} = \begin{pmatrix}
  \cos{\frac{i\pi}{6}} & -\sin{\frac{i\pi}{6}}\\ 
  \sin{\frac{i\pi}{6}} & \cos{\frac{i\pi}{6}}
\end{pmatrix} \begin{pmatrix}
  \cos({\frac{(k-4)\pi}{6}}) \\ 
  \sin({\frac{(k-4)\pi}{6}})
\end{pmatrix} = \begin{pmatrix}
  \cos({\frac{i\pi}{6}})\cos({\frac{(k-4)\pi}{6}}) - \sin({\frac{i\pi}{6}})\sin({\frac{(k-4)\pi}{6}})\\ 
  \sin({\frac{i\pi}{6}})\sin({\frac{(k-4)\pi}{6}}) + \cos({\frac{i\pi}{6}})\cos({\frac{(k-4)\pi}{6}})
\end{pmatrix} = \begin{pmatrix}
  \cos({\frac{(i+k-4)\pi}{6}}) \\ 
  \sin({\frac{(i+k-4)\pi}{6}})
\end{pmatrix}$

\end{center}
\end{block}
\end{frame}
\begin{frame}
\begin{block}{Rotating $v_3$}
\begin{center}
$R_{\frac{\pi}{6}}v_{3} = \begin{pmatrix}
  \cos{\frac{i\pi}{6}} & -\sin{\frac{i\pi}{6}}\\ 
  \sin{\frac{i\pi}{6}} & \cos{\frac{i\pi}{6}}
\end{pmatrix} \begin{pmatrix}
  \cos({\frac{(k-7)\pi}{6}}) \\ 
  \sin({\frac{(k-7)\pi}{6}})
\end{pmatrix} = \begin{pmatrix}
  \cos({\frac{i\pi}{6}})\cos({\frac{(k-7)\pi}{6}}) - \sin({\frac{i\pi}{6}})\sin({\frac{(k-7)\pi}{6}})\\ 
  \sin({\frac{i\pi}{6}})\sin({\frac{(k-7)\pi}{6}}) + \cos({\frac{i\pi}{6}})\cos({\frac{(k-7)\pi}{6}})
\end{pmatrix} = \begin{pmatrix}
  \cos({\frac{(i+k-7)\pi}{6}}) \\ 
  \sin({\frac{(i+k-7)\pi}{6}})
\end{pmatrix}$

\end{center}
\end{block}
\end{frame}

